\section{Arithmetic cancellation and the geometric remainder}
\label{sec:arithmeticgeometry}

\subsection{A transfer estimate with both inputs explicit}
\label{subsec:primitiveRHtransfer}

Let $B\subset\mathbb R^2$ be a fixed convex body containing the origin in
its interior, and normalize its gauge so that the shortest nonzero
integer vector has gauge one.  Put
\[
 A_B(u)=\#\bigl(uB\cap(\mathbb Z^2\setminus\{0\})\bigr),\qquad
 \Delta_B(u)=A_B(u)-a_Bu^2,\qquad a_B=\operatorname{Area}(B),
\]
and write
\[
 P_B(L)=\#\bigl(LB\cap\mathbb Z^2_{\rm prim}\bigr),\qquad
 E_B(L)=P_B(L)-\frac{a_B}{\zeta(2)}L^2.
\]
All primitive vectors are counted with their signs; an unoriented count
for a centrally symmetric body is half this count.
The notation $\Delta_B$ here is the all-point discrepancy $D_B$ of
Section~\ref{subsec:mobiusmodulus}.

There are two inputs in the primitive problem: the geometric discrepancy
$\Delta_B$, and cancellation in the M\"obius function.  A quantitative
bound for the second input does not by itself specify the first.  The
following elementary estimate keeps their roles separate.

\begin{proposition}[A large-divisor transfer estimate]
\label{prop:primitiveRHtransfer}
Suppose
\[
 M(x):=\sum_{n\le x}\mu(n),\qquad |M(x)|\le Cx\eta(x),
\]
where $\eta$ is nonnegative and nonincreasing on $[1,\infty)$.
Then, for every integer $1\le D\le L$,
\begin{equation}\label{eq:primitiveRHmaster}
 |E_B(L)|\ll_{B,C}
       \sum_{d\le D}|\Delta_B(L/d)|
       +\frac{L^2}{D}\eta(D).
\end{equation}
In particular, $M(x)\ll x^\theta$, with $0\le\theta<1$, gives
\begin{equation}\label{eq:primitiveRHpowerM}
 |E_B(L)|\ll_{B,\theta}
       \sum_{d\le D}|\Delta_B(L/d)|+L^2D^{\theta-2}.
\end{equation}
\end{proposition}

\begin{proof}
M\"obius inversion gives
\[
 P_B(L)=\sum_{d\le L}\mu(d)A_B(L/d).
\]
Set $U=L/D$.  Reversing the order of summation in its large-divisor
part gives the exact identity
\[
 \sum_{d>D}\mu(d)A_B(L/d)
   =\sum_{\substack{m\in\mathbb Z^2\setminus\{0\}\\\|m\|_B\le U}}
       \left(M\!\left(\frac{L}{\|m\|_B}\right)-M(D)\right).
\]
The summands with $\lfloor L/\|m\|_B\rfloor=D$ are zero.
Since $A_B(u)\ll_B u^2$ for $u\ge1$, partial summation gives
\[
 \sum_{\substack{m\ne0\\\|m\|_B\le U}}\|m\|_B^{-1}\ll_B U.
\]
Monotonicity of $\eta$ therefore bounds the large-divisor part by
\[
 \eta(D)\left(L\sum_{0<\|m\|_B\le U}\|m\|_B^{-1}
                      +D A_B(U)\right)
       \ll_B\frac{L^2}{D}\eta(D).
\]
The area term has the same tail bound, because partial summation yields
\[
 \left|\sum_{d>D}\frac{\mu(d)}{d^2}\right|
       \ll\frac{\eta(D)}{D}.
\]
Subtract the area main term in the small-divisor sum and bound its
remaining terms absolutely.  This proves the first estimate; take
$\eta(x)=x^{\theta-1}$ for the second.
\end{proof}

\subsection{An unconditional sublinear consequence of the Dini condition}
Define the nonincreasing geometric modulus
\[
 \omega_B(u)=\sup_{v\ge u}\frac{|\Delta_B(v)|}{v}.
\]
For $L/(D+1)\ge1$, direct comparison on the intervals
$[L/(d+1),L/d]$ gives
\begin{equation}\label{eq:primitiveRHmodulus}
 \sum_{d\le D}|\Delta_B(L/d)|
   \le 2L\int_{L/(D+1)}^L\omega_B(u)\frac{du}{u}.
\end{equation}
Consequently,
\begin{equation}\label{eq:primitiveDiniPNT}
 \int_1^\infty\omega_B(u)\frac{du}{u}<\infty
 \quad\Longrightarrow\quad E_B(L)=o_B(L).
\end{equation}
Here the arithmetic input is only the prime number theorem, in the
equivalent form $M(x)=o(x)$, not RH.  Indeed, use the nonincreasing
envelope $\eta(x)=\sup_{y\ge x}|M(y)|/y$, fix $U\ge2$, and choose
$D=\lfloor L/U\rfloor$.  After division by $L$, the large-divisor
bound tends to zero for this fixed $U$, whereas the small-divisor bound
is at most
$2\int_{U/2}^{\infty}\omega_B(u)\,du/u$ for all sufficiently large $L$.
Let $L\to\infty$ and then $U\to\infty$.
This strengthens the earlier absolute-value Dini argument, which gave
only $O_B(L)$.


\subsection{Fixed zero-free half-planes and the recent quasi-RH claims}
\label{subsec:quasiRHapplication}

For the Riemann zeta function and $1/2\le\theta_0<1$, the standard
zero-free-half-plane formulation is equivalent to
\begin{equation}
 \zeta(s)\ne0\quad(\Re s>\theta_0)
 \quad\Longleftrightarrow\quad
 M(x)=O_\varepsilon(x^{\theta_0+\varepsilon})
 \quad\hbox{for every }\varepsilon>0.
 \label{eq:zeroFreeMertens}
\end{equation}
Here nonvanishing excludes the pole at $s=1$, where $1/\zeta$
instead has a zero. We give the analytic step because a fixed zero-free
half-plane and a height-dependent zero-free region are different inputs.

Suppose first that the left side of~\eqref{eq:zeroFreeMertens} holds.
Fix $0<\delta<(1-\theta_0)/2$. At sufficiently large $|t|$, choose the
branch of $\log\zeta$ on the disk centered at $2+it$ of radius
$R_2=2-\theta_0-\delta/4$ that agrees with the Euler-product logarithm
at its center. This disk contains neither zeros nor the pole.
Polynomial growth of $\zeta$ in a fixed vertical strip gives
$\Re\log\zeta\ll_\delta\log(2+|t|)$ there.
Borel--Carath\'eodory bounds $|\log\zeta|$ by the same order on the
smaller disk of radius $R=2-\theta_0-\delta/2$. On the disk of radius
$r_0=1/2$ the absolutely convergent Euler-product logarithm is bounded.
The three-circles theorem, with $r=2-\theta_0-\delta<R$, now gives
\[
 |\log\zeta(\sigma+it)|\ll_\delta(\log(2+|t|))^{\lambda_\delta},
 \qquad \theta_0+\delta\le\sigma\le2,\qquad
 \lambda_\delta=\frac{\log(r/r_0)}{\log(R/r_0)}<1.
\]
Consequently, for every $\varepsilon>0$,
$1/\zeta(\sigma+it)\ll_{\delta,\varepsilon}(1+|t|)^\varepsilon$
on $\sigma\ge\theta_0+\delta$; bounded heights and $\sigma\ge2$
are handled by analyticity of the reciprocal and the Euler product.

For a half-integer $x$, truncated Perron inversion on
$c=1+1/\log x$, at height $T=x^2$, has error
$O(x\log^2x/T)$ because $|\mu(n)|\le1$. Move the contour to
$\sigma_0=\theta_0+\delta$. There are no poles of the reciprocal in
the intervening strip. The new vertical integral is
$O_{\delta,\varepsilon}(x^{\sigma_0}T^\varepsilon)$ and the two
horizontal integrals are $O_{\delta,\varepsilon}(xT^{\varepsilon-1})$.
Choosing $\delta$ and $\varepsilon$ sufficiently small gives every
exponent $\theta_0+\varepsilon'$ requested on the right side
of~\eqref{eq:zeroFreeMertens}. Half-integers suffice since $M$ is
constant between consecutive integers.

Conversely, for each compact subset of $\Re s>\theta_0$, choose
$\varepsilon>0$ smaller than its distance from the boundary. The bound
for $M$ makes $s\int_1^\infty M(x)x^{-s-1}\,dx$ locally uniformly
convergent and holomorphic there. Partial summation identifies it with
$1/\zeta(s)$ on $\Re s>1$, so analytic continuation precludes a zero
of $\zeta$ in $\Re s>\theta_0$. This equivalence uses the Euler product,
meromorphic continuation and vertical growth of $\zeta$; zero-freeness
of an arbitrary Dirichlet series would not alone suffice.

The OpenAI manuscript dated 30 September 2026 claims the half-plane
$\Re s>7/8$ for all finite-order Hecke $L$-functions over
$\mathbb Q(\sqrt{-3})$, and hence for Dirichlet $L$-functions,
including $\zeta$~\cite[Theorem~1.1]{OpenAIQuasiRHSeptember}.
The manuscript dated 5 October 2026 gives a separate proof claiming
$\Re s>11/12$ and describes it as a simpler intermediate result toward
the September theorem~\cite[Theorem~1.1 and the following paragraph]{OpenAIQuasiRHOctober}.
For both papers we use the theorem statements in the TeX sources at
repository revision \texttt{adc7f1241b42} (the initial release), as pinned
in the claim ledger and bibliography. We checked those statements,
not either complete proof or its formalization. All applications here are conditional on
the explicit arithmetic hypothesis in Proposition~\ref{prop:primitiveRHtransfer}.
The corresponding Mertens exponents would be $7/8+\varepsilon$,
$11/12+\varepsilon$, and, under full RH, $1/2+\varepsilon$.
The generalized assertions about other $L$-functions are unnecessary
for this primitive planar count.

\subsection{Logarithmic geometric savings}
The next comparison is conditional on an additional estimate for the
all-point discrepancy:
\begin{equation}\label{eq:primitiveRHgeometry}
 |\Delta_B(u)|\ll_B\frac{u}{(\log u)^\beta},\qquad u\ge3,
 \qquad\beta>0.
\end{equation}
No such estimate is asserted here for a McShane--Rivin norm ball.

First suppose $M(x)\ll x^\theta$ for some fixed $\theta<1$.
Choose any fixed $d$ with
$1/(2-\theta)<d<1$, and take $D=\lfloor L^d\rfloor$.
The large-divisor contribution in~\eqref{eq:primitiveRHpowerM} is
$O(L^{2-d(2-\theta)})$, with exponent strictly smaller than one.
The small-divisor contribution is
\[
 \ll L\int_{L^{1-d}/2}^L\frac{du}{u(\log u)^\beta}
 \ll_{\beta,d} L(\log L)^{1-\beta}.
\]
Thus
\begin{equation}\label{eq:primitiveRHlogresult}
 M(x)\ll x^\theta\ (\theta<1),\quad
 \Delta_B(u)\ll u(\log u)^{-\beta}
 \quad\Longrightarrow\quad
 E_B(L)\ll L(\log L)^{1-\beta}.
\end{equation}
In particular, the geometric threshold $\beta=1$ gives $O_B(L)$.
Any fixed power saving for $M$ suffices for this conclusion: full RH
does not improve the logarithmic exponent in this elementary argument.

For comparison, the classical unconditional Vinogradov--Korobov estimate,
recorded for example in Zhai~\cite[eq.~(1.5)]{ZhaiPrimitive2003}, has the form
\[
 M(x)\ll x\exp\!\left(-c(\log x)^{3/5}
                              (\log\log x)^{-1/5}\right).
\]
Write $F(L)=(\log L)^{3/5}(\log\log L)^{-1/5}$ and take
$D=\lfloor L\exp(-aF(L))\rfloor$ with a sufficiently small fixed
$a>0$.  Since $F(D)\sim F(L)$, the large-divisor part is
$O(L\exp(-c'F(L)))$.  The small-divisor part is bounded by
\[
 L\int_{\exp(aF(L))/2}^L\frac{du}{u(\log u)^\beta}.
\]
The resulting elementary comparison is
\[
\begin{array}{c|c|c}
 \text{geometric exponent }\beta
 &\text{unconditional VK input}
 &M(x)\ll x^\theta,\ \theta<1\\ \hline
 0<\beta<1
 &O\!\left(L(\log L)^{1-\beta}\right)
 &O\!\left(L(\log L)^{1-\beta}\right)\\[2pt]
 \beta=1
 &O(L\log\log L)
 &O(L)\\[2pt]
 \beta>1
 &O\!\left(LF(L)^{1-\beta}\right)
 &O\!\left(L(\log L)^{1-\beta}\right).
\end{array}
\]
These are consequences of one elementary splitting argument, not
assertions of optimality.  The constants can depend on $B$, $\beta$,
the arithmetic exponent, and the chosen cutoff exponent.

\subsection{Power-saving geometric discrepancies and established benchmarks}
If instead $|\Delta_B(u)|\ll u^\alpha$ with $0\le\alpha<1$, then
\[
 \sum_{d\le D}|\Delta_B(L/d)|\ll L^\alpha D^{1-\alpha}.
\]
Balancing this with $L^2D^{\theta-2}$ gives
\begin{equation}\label{eq:primitiveRHbalancedpower}
 D\asymp L^{(2-\alpha)/(3-\alpha-\theta)},\qquad
 E_B(L)\ll L^{\kappa(\alpha,\theta)},\qquad
 \kappa(\alpha,\theta)=\frac{2-\alpha\theta}{3-\alpha-\theta}.
\end{equation}
For example, $\alpha=2/3$ and a Mertens exponent arbitrarily close
to $1/2$ give $O_\varepsilon(L^{10/11+\varepsilon})$.
Exponents arbitrarily close to $7/8$ and $11/12$ give, respectively,
$O_\varepsilon(L^{34/35+\varepsilon})$ and
$O_\varepsilon(L^{50/51+\varepsilon})$ by this method.
These examples concern a hypothetical or separately established
geometric estimate with $\alpha=2/3$; they do not assign that exponent
to a McShane--Rivin ball.  Stronger results for smooth curved domains
use additional structure and should not be confused with this general
transfer bound.

With the same power-saving geometric input, the unconditional VK
estimate already gives
\[
 E_B(L)\ll L\exp(-c_B F(L))
\]
for some $c_B>0$, by choosing
$D=L\exp(-aF(L))$ and balancing the two exponential factors.
Thus sublinear primitive errors in suitable geometries do not require RH.


This unconditional sublinear conclusion is a classical one for
sufficiently smooth, positively curved domains: Zhai records it in
equation~(1.4) of~\cite{ZhaiPrimitive2003}. His parameter is the area
scale $x=L^2$, so his $O(x^{1/2}\exp(-cF(x)))$ becomes
$O(L\exp(-c'F(L)))$ here. The hypotheses include finite nonzero
curvature throughout and a $C^4$ Gauss map. Those hypotheses do not
hold for a McShane--Rivin ball.

Under RH, Baker~\cite[Theorem~1]{BakerPrimitive2010} proves
$O_\varepsilon(L^{10/13+\varepsilon})$ for a specified class of
piecewise $C^4$ domains whose arcs have nonzero curvature, with
additional $C^4$ reciprocal-curve regularity and rational endpoint
tangents, as specified by hypotheses (H1)--(H4) on pages~267--268.
This is an established comparison result, not a claim about the best
current exponent or about the flat norm balls considered here.
In the opposite direction, the unconditional polygon obstruction
$\Omega_\pm(L\sqrt{\log\log L})$ recalled in
Section~\ref{subsec:rationalpolygonbenchmark} persists regardless of RH.
It belongs to rational polygons and is not a lower bound for strictly
convex bodies. These examples show why the shape hypothesis must
accompany any arithmetic hypothesis.

\subsection{The remaining geometric obstruction}
Strict convexity supplies $\Delta_B(u)=o(u)$ and hence
$\omega_B(u)\to0$, but no specified rate for that convergence.
Under any fixed arithmetic exponent $\theta<1$, the preceding
power cutoff yields a bound of the form
\[
 |E_B(L)|\ll L\log L\,\omega_B(cL^{1-d})
                      +L^{2-d(2-\theta)}.
\]
It is again $o(L\log L)$, without a universal logarithmic saving
specified by this argument.  This is a limitation of the stated inputs
and this transfer estimate, not an impossibility theorem for more
special convex bodies.

One cannot simply multiply a bound for $M$ by the supremum of the
weights $\Delta_B(L/d)$.  Abel summation also sees their variation.
On a dyadic interval $d\asymp D$, the counting part of these weights
can have total variation of order $(L/D)^2$, despite a much smaller
pointwise discrepancy.  The large-divisor estimate above pays for
that variation explicitly.  Further cancellation would require
additional information about the geometric weights or their correlation
with $\mu(d)$.

Equivalently, for $\Re s>2$ the gauge zeta functions satisfy
\[
 Z_B(s)=\sum_{m\ne0}\|m\|_B^{-s},\qquad
 Z_B^{\rm prim}(s)=\frac{Z_B(s)}{\zeta(s)}.
\]
The estimate $\Delta_B(u)=o(u)$ provides the representation
\[
 Z_B(s)=\frac{a_Bs}{s-2}
          +s\int_1^\infty\Delta_B(u)u^{-s-1}\,du,
 \qquad \Re s>1.
\]
A zero-free result for the denominator does not on its own continue
the geometric numerator through the line $\Re s=1$, nor provide the
vertical growth bounds needed for a sharp counting inversion.  This is
the analytic reason to keep the arithmetic and geometric inputs separate.


For a body satisfying the infinite-flatness and rational-normal hypotheses
of Proposition~\ref{prop:normballflatrow}, the prime-row construction gives
$\Omega(L^{1-\varepsilon})$ along an unbounded sequence for every
$\varepsilon>0$. This excludes a fixed sublinear power bound for that
body even if RH holds. It does not exclude $O(L)$, or even $o(L)$:
the quantitative lower bound is $L W(h/L^2)$ up to constants, and
$W(h/L^2)\to0$ under strict convexity. A useful next geometric target is therefore a quantitative modulus.
One possible target is $\Delta_B(u)=O(u/\log u)$, provided it is
compatible with that body's local widths. Under the same supporting-row,
two-sided-graph and irrational-contact hypotheses, the row is empty at
$L=n/h$ and contributes $w_n(L^{-1})+O(1)$ all points at $L+L^{-1}$.
The all-point main term changes by only $a_B(2+L^{-2})$, so
\[
 \max\{|\Delta_B(L)|,|\Delta_B(L+L^{-1})|\}
 \ge \frac L2W(h/L^2)-O_B(1).
\]
There is no totient density factor in this all-point estimate, and it
holds for every sufficiently large integer $n$, without primality.
Consequently the proposed target requires $W(h/L^2)=O(1/\log L)$
along these contact scales. Strict
convexity and infinite flatness alone do not give that condition;
the model $f(s)=\exp[-\exp(|s|^{-\beta})]$ in
Section~\ref{subsec:flatnessmodels} violates it. Establishing a compatible
modulus for any specified McShane--Rivin ball remains a separate problem.

The all-point target has a numerical diagnostic independent of the area
constant, just as in Section~\ref{subsec:dyadicdiagnostic}. Define
\[
 \mathcal D_A(L)=A_B(2L)-4A_B(L).
\]
The dyadic inverse identity proves the equivalence
\begin{equation}
 \Delta_B(L)=O(L/\log L)
 \quad\Longleftrightarrow\quad
 \mathcal D_A(L)=O(L/\log L).
 \label{eq:allPointLogDiagnostic}
\end{equation}
Indeed, the reverse implication follows from
$\Delta_B(L)=-\sum_{j\ge0}4^{-j-1}\mathcal D_A(2^jL)$
and $\log(2^jL)\ge\log L$. For a symmetric torus norm ball,
let $\lambda_\gamma$ be the simple-geodesic length in the same
normalization as the gauge of $B$. The length list already used in
Section~\ref{subsec:torustraceexperiment}, rescaled when imposing the
unit-shortest-vector convention, gives the exact identity
\[
 A_B(L)=2\sum_{\substack{\gamma\text{ unoriented simple}\\\lambda_\gamma\le L}}
               \left\lfloor\frac L{\lambda_\gamma}\right\rfloor.
\]
Thus the proposed geometric threshold can be investigated without
estimating the norm-ball area or introducing a new enumeration of
curve types. No finite numerical range proves the uniform bound
in~\eqref{eq:allPointLogDiagnostic}.

The full campaign \path{experiments/torus_campaign/run.py} implements
this signed-multiple diagnostic at its dyadic cutoffs, with the all-point
rows and differences in \path{artifacts/torus_campaign/result.json}.
The floor comparisons use exact rational representations of the stored
binary64 lengths. Membership for the true geodesic lengths remains
numerical; the exact identity does not certify the required length or
completeness assumptions. The campaign uses hyperbolic lengths
$\ell_\gamma$ without rescaling. If
$s_X$ is the shortest hyperbolic length, the unit-shortest-vector gauge
uses $\lambda_\gamma=\ell_\gamma/s_X$ and hence
$A_B(L)=A_{\rm hyp}(s_XL)$. The campaign's hyperbolic dyadic rows must
be interpreted with that change of units; they do not establish a
uniform logarithmic bound. The separate finite companion
\path{experiments/counting_mixing/run.py} checks rational examples of
M\"obius inversion, signed multiples, row selectors, finite dyadic
identities and conditional balance exponents. These finite checks do
not supply either the required geometric modulus or an arithmetic
hypothesis. Source, independent-review and generation records are in
\path{docs/execution/counting-mixing-review-2026-11-01/}.

